\section{Hecke 代数在 Jacobian 上的作用}
\label{sec:hecke-on-jacobian}
% ============================================================================

到上一节为止，我们已经把 $J_0(N)$ 构造成一个由权 $2$ 模形式和周期格组成的 Abel 簇。
但这还只解决了“几何容器”的问题，尚未解决“算术结构”的问题。
在第\ref{ch:hecke}章里，Hecke 算子只是 $\Sk_2(\GammaO(N))$ 上的一族交换线性算子；
如果它们不能进一步作用在 $J_0(N)$ 上，那么我们仍然无法把特征形式切成真正的几何因子。

\textbf{我们遇到了什么问题？}
单看向量空间 $\Sk_2(\GammaO(N))$，特征形式分解只是线性代数。
而模性问题真正需要的是 Abel 簇之间的同态、商与同种。
所以我们必须回答：Hecke 算子不仅能作用在微分上，能否也作用在 Jacobian 本身上？

\textbf{这个概念的核心思想是什么？}
答案来自\textbf{对应（correspondence）}\index{对应!correspondence}。
Hecke 算子并不是凭空写下的一串公式；
它来自模曲线之间的一对有限映射
\[
  X_0(N) \xleftarrow{\ \alpha_n\ } X_0(N;n) \xrightarrow{\ \beta_n\ } X_0(N),
\]
其中 $X_0(N;n)$ 参数化带有一个 $N$ 级循环子群和一个次数 $n$ 循环同源的数据。
把这对映射先拉回再推出，就得到 Jacobian 上的自同态。

\textbf{引入之后能得到什么？}
一旦这些对应真的给出 $J_0(N)$ 的自同态，
Hecke 代数 $\mathbb T$ 就进入了 $\End(J_0(N))$。
这样一来，每个特征形式的特征值系统都会对应 Jacobian 中的某个几何因子，
这正是下一节构造 $A_f$ 的起点。

\begin{figure}[ht]
\centering
\begin{tikzpicture}[>=Stealth, font=\small, node distance=2.7cm]
  \node[draw, rounded corners, fill=blue!8, minimum width=3.4cm, minimum height=1cm] (x1) {$X_0(N)$};
  \node[draw, rounded corners, fill=green!8, right of=x1, minimum width=4.0cm, minimum height=1cm] (y) {$X_0(N;n)$};
  \node[draw, rounded corners, fill=blue!8, right of=y, minimum width=3.4cm, minimum height=1cm] (x2) {$X_0(N)$};

  \draw[->, thick] (y) -- node[above] {$\alpha_n$} (x1);
  \draw[->, thick] (y) -- node[above] {$\beta_n$} (x2);

  \node[draw, rounded corners, fill=orange!10, below=2.1cm of y, minimum width=6.0cm, minimum height=1cm] (j)
    {$J_0(N)$ 上的自同态 $T_n=(\beta_n)_*\alpha_n^*$};

  \draw[->, thick] ($(x1.south)+(0.1,0)$) -- ++(0,-0.85) -| (j.west);
  \draw[->, thick] ($(x2.south)+(-0.1,0)$) -- ++(0,-0.85) -| (j.east);

  \node[below=0.8cm of j] {在余切空间上，$H^0(J_0(N),\Omega^1)\cong \Sk_2(\GammaO(N))$，且 $T_n$ 的作用正是经典 Hecke 算子。};
\end{tikzpicture}
\caption{Hecke 算子来自模曲线之间的对应：先把数据拉回到 $X_0(N;n)$，再推出回 $X_0(N)$，于是得到 Jacobian 上的自同态。}
\label{fig:ch06-hecke-jacobian}
\end{figure}


\begin{definition}[Hecke 对应]
\label{def:hecke-correspondence-jacobian}
对每个正整数 $n$，记 $X_0(N;n)$ 为参数化三元组 $(E,C,D)$ 的紧模曲线，
其中 $C\subset E$ 是阶 $N$ 的循环子群，$D\subset E$ 是阶 $n$ 的循环子群，且 $C\cap D=0$。
定义两个有限映射
\begin{align*}
  \alpha_n(E,C,D)&=(E,C),\\
  \beta_n(E,C,D)&=\bigl(E/D,(C+D)/D\bigr).
\end{align*}
它们组成 $X_0(N)$ 上的一个对应。
在除子类群 $\operatorname{Pic}^0(X_0(N))\cong J_0(N)$ 上，
相应的 Hecke 算子定义为
\[
  T_n=(\beta_n)_*\circ \alpha_n^*.
\]
\end{definition}

\begin{proposition}[Hecke 对应确实作用在 $J_0(N)$ 上]
\label{prop:hecke-correspondence-on-jacobian}
\cref{def:hecke-correspondence-jacobian} 中的 $T_n$ 把次数为 $0$ 的除子类送到次数为 $0$ 的除子类，
并且对主除子给出主除子。
因此它诱导 $J_0(N)$ 上的一个群同态，仍记为 $T_n$。
\end{proposition}

\begin{proof}
先看次数。
若 $D=\sum_P m_P[P]$ 是 $X_0(N)$ 上的除子，则
\[
  \deg \alpha_n^*D=(\deg \alpha_n)\deg D,
\]
因为每个点的原像总重数等于映射次数。
于是若 $\deg D=0$，则 $\deg \alpha_n^*D=0$。
再对 $Y=X_0(N;n)$ 上的除子 $E=\sum_Q n_Q[Q]$，有
\[
  \deg (\beta_n)_*E=\sum_Q n_Q=\deg E,
\]
因为推出只是把位于同一像点上的系数相加。
故 $T_n$ 保持次数 $0$。

再看主除子。
若 $D=(f)$ 是某个亚纯函数 $f$ 的除子，则
\[
  \alpha_n^*D=(f\circ \alpha_n),
\]
仍是主除子。
另一方面，若 $E=(g)$ 是 $Y$ 上的主除子，则其推出是
\[
  (\beta_n)_*E=\bigl(N_{\beta_n}(g)\bigr),
\]
其中 $N_{\beta_n}(g)$ 是沿有限映射 $\beta_n$ 的范数函数。
局部地说，$N_{\beta_n}(g)$ 是各原像点处函数值按分歧指数加权的乘积，
因此其零极点阶数恰好是 $E$ 推出后的系数。
所以推出也把主除子送到主除子。

综上，$T_n=(\beta_n)_*\alpha_n^*$ 在 $\operatorname{Pic}^0(X_0(N))$ 上良定义，
经 Abel 定理与 Jacobian 的识别，就得到 $J_0(N)$ 上的群同态。
\end{proof}

\begin{proposition}[Hecke 作用与权 $2$ 形式一致]
\label{prop:hecke-on-cotangent-jacobian}
把 $H^0(J_0(N),\Omega^1)$ 与 $\Sk_2(\GammaO(N))$ 通过\cref{prop:j0n-from-s2} 识别。
则 Jacobian 上的 Hecke 自同态 $T_n$ 在余切空间上的作用，正是第\ref{ch:hecke}章定义的经典 Hecke 算子。
\end{proposition}

\begin{proof}
对任意紧曲线上的对应
\[
  X \xleftarrow{\alpha} Y \xrightarrow{\beta} X,
\]
其在 $\operatorname{Pic}^0(X)$ 上的作用为 $(\beta)_*\alpha^*$。
对余切空间求微分，得到的就是微分形式上的拉回再推出：
\[
  H^0(X,\Omega^1) \xrightarrow{\ \alpha^*\ } H^0(Y,\Omega^1)
  \xrightarrow{\ \beta_*\ } H^0(X,\Omega^1).
\]
原因很直接：在 Jacobian 的原点附近，点的变化由积分微分控制，
而对应对除子做的“拉回--推出”，在线性化后正是对微分做“拉回--迹映射”。

对于模曲线的 Hecke 对应，这个复合映射恰是经典定义中的 Hecke 算子。
第\ref{ch:hecke}章中无论通过双陪集还是通过同源计数得到的 $T_n$，
本质上都来自同一个模解释：
把一个点 $(E,C)$ 拉回到所有带有阶 $n$ 子群 $D$ 的三元组 $(E,C,D)$，
再把它们推到各个商曲线 $(E/D,(C+D)/D)$。
因此在
\[
  H^0\bigl(X_0(N),\Omega^1\bigr)\cong \Sk_2\bigl(\GammaO(N)\bigr)
\]
上的作用与经典 Hecke 算子完全一致。
\end{proof}

\begin{definition}[Hecke 代数在 Jacobian 中的实现]
\label{def:hecke-algebra-jacobian}
记 $\mathbb T$ 为由所有 $T_n$ 在 $\End(J_0(N))$ 中生成的子环，
称为 $J_0(N)$ 上的\textbf{Hecke 代数（Hecke algebra）}\index{Hecke代数!Hecke algebra}。
\end{definition}

\begin{proposition}[Jacobain 上的作用由余切空间决定]
\label{prop:endomorphism-determined-by-cotangent}
设 $A$ 是复 Abel 簇，$u\in\End(A)$。
若 $u$ 在原点余切空间 $H^0(A,\Omega^1)$ 上的作用为零，
则 $u=0$。
\end{proposition}

\begin{proof}
设 $du_0:T_0A\to T_0A$ 是 $u$ 在原点的微分。
余切空间上的作用为零等价于对偶线性映射 $du_0^*$ 为零，因而 $du_0=0$。
对任意点 $P\in A$，利用群同态性质
\[
  u\circ \tau_P=\tau_{u(P)}\circ u,
\]
其中 $\tau_P$ 表示平移。
对上式在原点求微分，得到
\[
  du_P=d\tau_{u(P),0}\circ du_0\circ d\tau_{P,0}^{-1}=0.
\]
所以 $u$ 在每一点的微分都为零。
于是 $u$ 是连通复流形上的局部常值映射，只能是常值。
又因 $u$ 是群同态，必有 $u(0)=0$，故这个常值只能是零元。
因此 $u=0$。
\end{proof}

\begin{corollary}[Hecke 代数交换]
\label{cor:hecke-algebra-commutative-jacobian}
$\mathbb T\subset \End(J_0(N))$ 是交换环。
\end{corollary}

\begin{proof}
由\cref{prop:hecke-on-cotangent-jacobian}，
每个 $T_n$ 在余切空间上的作用就是第\ref{ch:hecke}章的经典 Hecke 算子。
这些算子彼此交换。
因此对任意 $m,n$，交换子
\[
  [T_m,T_n]=T_mT_n-T_nT_m
\]
在 $H^0(J_0(N),\Omega^1)$ 上的作用为零。
由\cref{prop:endomorphism-determined-by-cotangent}，该交换子本身只能是零。
故 $T_mT_n=T_nT_m$，于是 $\mathbb T$ 交换。
\end{proof}

\begin{proposition}[特征形式在 Jacobian 上留下特征值系统]
\label{prop:eigenform-gives-character}
若 $f\in \Sk_2(\GammaO(N))$ 是归一化 Hecke 特征形式，满足
\[
  T_nf=a_n(f)f \qquad (n\ge 1),
\]
则映射
\[
  \lambda_f:\mathbb T\longrightarrow \C,
  \qquad
  T_n\longmapsto a_n(f)
\]
是一个环同态。
\end{proposition}

\begin{proof}
因为 $f$ 同时是所有 $T_n$ 的特征向量，所以对任意多项式表达式 $P(T_1,\dots,T_r)$ 都有
\[
  P(T_1,\dots,T_r)f=P\bigl(a_1(f),\dots,a_r(f)\bigr)f.
\]
特别地，对任意 $S,T\in\mathbb T$，若记 $Sf=\lambda_f(S)f$、$Tf=\lambda_f(T)f$，则
\[
  (S+T)f=(\lambda_f(S)+\lambda_f(T))f,
  \qquad
  (ST)f=\lambda_f(S)\lambda_f(T)f.
\]
于是
\[
  \lambda_f(S+T)=\lambda_f(S)+\lambda_f(T),
  \qquad
  \lambda_f(ST)=\lambda_f(S)\lambda_f(T).
\]
故 $\lambda_f$ 是环同态。
\end{proof}

\textbf{这一节真正得到的东西。}
Hecke 算子已经不再只是模形式空间上的矩阵，
而是 Jacobian 上的几何自同态。
因此特征值系统也不再只是数字，而会决定一个真正的 Abel 簇商。
下一节的 $A_f$ 正是按这个思路定义出来的。

\begin{definition}[Rosati 对合]
\label{def:rosati-involution}
设 $A$ 是带极化 $\lambda:A\to A^{\vee}$ 的 Abel 簇。
对任意 $u\in \End(A)$，定义
\[
  u^{\dagger}=\lambda^{-1}\circ \widehat{u}\circ \lambda,
\]
其中 $\widehat{u}:A^{\vee}\to A^{\vee}$ 是对偶同态。
这给出 $\End(A)$ 上一个正定对合，称为\textbf{Rosati 对合}。
\end{definition}

\begin{proposition}[Hecke 算子对 Rosati 对合自伴]
\label{prop:hecke-selfadjoint-rosati}
对 $J_0(N)$ 的自然主极化，
每个 Hecke 算子都满足
\[
  T_n^{\dagger}=T_n.
\]
\end{proposition}

\begin{proof}
把 $T_n$ 看成对应
\[
  X_0(N) \xleftarrow{\ \alpha_n\ } X_0(N;n) \xrightarrow{\ \beta_n\ } X_0(N)
\]
诱导的自同态 $(\beta_n)_*\alpha_n^*$。
对 Jacobian 的主极化而言，Rosati 对合正是把一个对应换成它的转置对应，
所以
\[
  T_n^{\dagger}=(\alpha_n)_*(\beta_n)^*.
\]
而对 Hecke 对应，交换两边箭头只不过是把次数 $n$ 的循环同源换成其对偶同源；
这并不会改变对应本身。
因此转置对应与原对应自然同构，于是 $T_n^{\dagger}=T_n$。
\end{proof}

\begin{corollary}[实特征值与 Petersson 内积]
\label{cor:hecke-real-eigenvalues}
在识别
\[
  H^0\bigl(J_0(N),\Omega^1\bigr)\cong \Sk_2\bigl(\GammaO(N)\bigr)
\]
下，Rosati 对合诱导的 Hermite 型与第3章的 Petersson 内积只差一个正实常数。
因此 $T_n$ 在 $\Sk_2(\GammaO(N))$ 上是自伴 Hermitian 算子，
从而其特征值都是实数。
\end{corollary}

\begin{proof}[说明]
对 Jacobian 的自然主极化，Riemann 双线性关系给出的 Hermite 型与曲线上的积分配对一致；
在模曲线的情形，这个积分配对正是 Petersson 内积的几何化版本。
既然 \cref{prop:hecke-selfadjoint-rosati} 已说明 $T_n$ 关于 Rosati 对合自伴，
那么它关于 Petersson 内积也自伴。
有限维 Hermitian 空间上线性代数的谱定理立刻推出：
$T_n$ 的全部特征值都属于 $\R$。
\end{proof}
